\documentclass[a4paper,12pt]{article}
\usepackage[utf8]{inputenc}
\usepackage{geometry}
\usepackage{fancyhdr}
\usepackage{xcolor}
\usepackage{minted}
\usepackage{graphicx}
\usepackage{hyperref}
\usepackage{enumitem}

% Geometry setup
\geometry{top=2.5cm, bottom=2.5cm, left=2.5cm, right=2.5cm}

% Header and Footer
\pagestyle{fancy}
\fancyhf{}
\lhead{\textbf{CMP9134: Software Engineering}}
\rhead{\textbf{Week 6}}
\lfoot{University of Lincoln}
\rfoot{Page \thepage}

% Color definitions
\definecolor{lightgray}{rgb}{0.95,0.95,0.95}

\begin{document}

%----------------------------------------------------------------------------------------
%   TITLE SECTION
%----------------------------------------------------------------------------------------

\begin{center}
    \LARGE \textbf{Lab Sheet 6: HCI \& Design Thinking} \\
    \vspace{0.5cm}
    \large \textbf{Objectives}
\end{center}

\noindent By completing today's workshop, you will:
\begin{itemize}
    \item Apply Human-Computer Interaction (HCI) principles to design a User Interface (UI).
    \item Create a low-fidelity wireframe prototype of your Ground Control Station dashboard.
    \item Conduct a Heuristic Evaluation of your prototype with human peers and an AI assistant.
    \item Build an interactive "Limited functionality simulation" using frontend code or Generative AI tools.
    \item Perform an automated Accessibility audit to ensure LEPSI/WCAG compliance.
\end{itemize}

\vspace{0.5cm}
\hrule
\vspace{0.5cm}

\section*{Introduction: Prototyping Your Dashboard}
In your Assessment 1 project, you are tasked with building a web-based Ground Control Station to interact with a Virtual Robot. A common mistake in software engineering is jumping straight into writing HTML, CSS, and JavaScript before knowing what the user actually needs. Today, we will use the \textbf{Design Thinking} methodology to Prototype and Test your interface before you write any production code.

%----------------------------------------------------------------------------------------
%   TASK 1
%----------------------------------------------------------------------------------------

\section*{Task 1: Wireframing the Web Dashboard}

\textbf{Goal:} Design the layout of your Ground Control Station and map out one specific feature in detail.

\textbf{The Concept:} Wireframes are low-fidelity visual representations of your interface. They focus on layout, information architecture, and user flow rather than colors and typography.

\vspace{0.3cm}
\textbf{Action:}
\begin{enumerate}
    \item Grab a pen and paper, or open a digital wireframing tool like \href{https://www.figma.com/}{Figma} or \href{https://excalidraw.com/}{Excalidraw}.
    \item \textbf{The Main Dashboard:} Sketch the primary view of your Ground Control Station. Ensure you include:
    \begin{itemize}
        \item A navigation bar/menu.
        \item A display for the Robot's current status (Battery, Status, X/Y Position).
        \item A visual indication of the current user's Role (Commander, Viewer, etc.).
    \end{itemize}
    \item \textbf{Deep-Dive Feature:} Pick one complex feature (e.g., The "Move Robot" control panel or the "Mission Audit Log") and draw a detailed, zoomed-in sequence of how the user interacts with it.
    \item \textit{Check:} Did you apply Fitts's Law? Are your primary action buttons (like "Move" or "Emergency Stop") large and easily accessible?
\end{enumerate}

%----------------------------------------------------------------------------------------
%   TASK 2
%----------------------------------------------------------------------------------------

\section*{Task 2: Peer and AI Usability Evaluation}

\textbf{Goal:} Evaluate the effectiveness of your design using established heuristics before building it.

\textbf{The Concept:} We will simulate a Heuristic Evaluation. You will gather feedback from two different sources: your human peers (who understand the project context) and a multimodal AI model.

\vspace{0.3cm}
\textbf{Action:}
\begin{enumerate}
    \item \textbf{Peer Evaluation:} Show your sketches to at least \textbf{two other students} in the room. Ask them to evaluate your UI against Shneiderman's Golden Rules (e.g., "Is there informative feedback when I click 'Move'?", "Is there an easy reversal/undo option?"). Write down their feedback.
    \item \textbf{AI Evaluation:} Take a clear photo or screenshot of your wireframe.
    \item Open a multimodal AI (with your student account, you have free access to GitHub Copilot Pro: \href{https://github.com/copilot}{https://github.com/copilot})  and upload your image.
    \item Use the following prompt to ask the AI for a formal critique:
    \begin{center}
    \fbox{%
    \begin{minipage}{0.9\linewidth}
\textit{
You are an expert UX/UI Researcher. I have uploaded a low-fidelity wireframe of a Ground Control Station dashboard used to pilot a Virtual Robot. 
Please conduct a Heuristic Evaluation of this interface based on Norman's 7 Principles and Shneiderman's 8 Golden Rules. 
Identify at least 3 usability issues and suggest specific improvements.
}
    \end{minipage}%
    }
    \end{center}
    \item \textbf{Refine:} Compare the feedback from your peers and the AI. Redraw or update your wireframe to fix the identified usability issues.
\end{enumerate}

%----------------------------------------------------------------------------------------
%   TASK 3
%----------------------------------------------------------------------------------------

\section*{Task 3: Limited Functionality Simulation (Interactive Prototype)}

\textbf{Goal:} Create a working, interactive prototype of your dashboard in the browser.

\textbf{The Concept:} In HCI, a "Limited functionality simulation" is a prototype that looks and feels like the final product but doesn't have a real backend connected to it. You must convert your refined wireframe from Task 2 into actual frontend code (HTML/CSS/JS). 

\vspace{0.3cm}
\textbf{Action (Choose Path A or Path B):}

\textbf{$\rightarrow$Path A: For students with Web Development Experience}
\begin{itemize}
    \item Manually code your dashboard based on your wireframe.
    \item Instead of writing all the CSS from scratch, we highly recommend using a \textbf{UI Framework}. Frameworks provide pre-built, standardized, and accessible components (buttons, navbars, grid systems).
    \item \textbf{Recommended Frameworks:} \href{https://getbootstrap.com/}{Bootstrap}, \href{https://tailwindcss.com/}{Tailwind CSS}, or your favourite one.
\end{itemize}

\textbf{$\rightarrow$Path B: For beginners (AI-Assisted)}
\begin{itemize}
    \item If you have zero frontend experience, you can use Generative AI tools (with your student account, you have free access to GitHub Copilot Pro: \href{https://github.com/copilot}{https://github.com/copilot}) to translate your sketch into code.
    \item Upload your final wireframe and provide a prompt describing the styling you want, such as:
    \begin{center}
    \fbox{%
    \begin{minipage}{0.9\linewidth}
\textit{Act as an expert Frontend Developer. Look at the wireframe I uploaded for a Robot Ground Control Station. 
Please generate a single-file "Limited functionality simulation" (HTML, CSS, and vanilla JavaScript) of this dashboard using the Bootstrap 5 framework. 
Make the styling modern, using a dark-mode theme suitable for a robotics control room. 
Make the "Move" button interactive so that clicking it shows a mock success notification.}
    \end{minipage}%
    }
    \end{center}
    \item Of course, the AI-generated code will not be perfect. You will need to review it, understand it, and make manual adjustments to ensure it matches your design and is functional. \textbf{Importantly, as with any AI usage, document your process as detailed in the assessment brief.}
\end{itemize}

\textbf{Final Step for both paths:} Save the resulting code into your project repository and open it in your web browser to test. This will serve as the foundational starting point for your Assessment 1 frontend!

%----------------------------------------------------------------------------------------
%   TASK 4 (STRETCH)
%----------------------------------------------------------------------------------------
%----------------------------------------------------------------------------------------
%   TASK 4 (STRETCH)
%----------------------------------------------------------------------------------------

\section*{Task 4: Accessibility Audit}

\textbf{Goal:} Ensure your interactive prototype is compliant with LEPSI and WCAG accessibility standards without relying on specific browser extensions.

\textbf{The Concept:} A beautiful interface is useless if visually impaired or motor-impaired users cannot navigate it. We will conduct a three-part browser-agnostic accessibility audit focusing on Structure, Contrast, and Keyboard Navigation.

\vspace{0.3cm}
\textbf{Action:}
\begin{enumerate}
    \item \textbf{Structural Audit (W3C):} Go to the \href{https://validator.w3.org/#validate_by_upload}{W3C Markup Validation Service}. Select "Validate by File Upload" and upload your html file. 
    \begin{itemize}
        \item Look for structural accessibility errors, such as missing \texttt{alt} attributes on images, missing \texttt{<label>} tags for inputs, or missing \texttt{<title>} tags.
    \end{itemize}
    \item \textbf{Color Contrast Audit (WebAIM):} Open the \href{https://webaim.org/resources/contrastchecker/}{WebAIM Contrast Checker} in your browser. 
    \begin{itemize}
        \item Enter the exact HEX color codes of your dark-mode background and your text. 
        \item Does it pass the WCAG AA standard (a ratio of at least 4.5:1)? If it fails, adjust your CSS colors until it passes.
    \end{itemize}
    \item \textbf{Manual Keyboard Audit:} Open your HTML file in \textit{any} web browser. Take your hand off your mouse/trackpad.
    \begin{itemize}
        \item Try to navigate to your "Move Robot" button using \textbf{ONLY} the \texttt{Tab} key, and "click" it using the \texttt{Enter} key. 
        \item Is it visually obvious which element currently has keyboard focus? If not, you need to add \texttt{:focus} styles to your CSS!
    \end{itemize}
    \item \textbf{Fix the Code:} Update your HTML/CSS to resolve any issues discovered during these three tests.
\end{enumerate}

%----------------------------------------------------------------------------------------
%   RESOURCES
%----------------------------------------------------------------------------------------

\section*{Further Resources}

\begin{itemize}
    \item \textbf{Bootstrap Documentation:} \href{https://getbootstrap.com/docs/5.3/getting-started/introduction/}{https://getbootstrap.com/} - \textit{The most popular UI framework for rapidly building responsive layouts.}
    % \item \textbf{v0 by Vercel:} \href{https://v0.dev/}{https://v0.dev/} - \textit{An AI tool specifically designed to generate React/HTML UIs from text prompts and images.}
    \item \textbf{W3C Markup Validator:} \href{https://validator.w3.org/}{https://validator.w3.org/} - \textit{The official tool to check your HTML for structural and accessibility errors.}
    \item \textbf{WebAIM Contrast Checker:} \href{https://webaim.org/resources/contrastchecker/}{https://webaim.org/resources/contrastchecker/} - \textit{Ensure your UI colors are legally compliant with WCAG guidelines.}
    \item \textbf{10 Usability Heuristics (NN/g):} \href{https://www.nngroup.com/articles/ten-usability-heuristics/}{nngroup.com/articles/ten-usability-heuristics/}
\end{itemize}

\end{document}

%----------------------------------------------------------------------------------------
%   RESOURCES
%----------------------------------------------------------------------------------------

% \section*{Further Resources}

% \begin{itemize}
%     \item \textbf{Bootstrap 5 Documentation:} \href{https://getbootstrap.com/docs/5.3/getting-started/introduction/}{https://getbootstrap.com/} - \textit{One of the most popular UI framework for rapidly building responsive layouts.}
%     \item \textbf{Google Lighthouse (Accessibility Auditing):}\\ \href{https://developer.chrome.com/docs/lighthouse/accessibility/}{developer.chrome.com/docs/lighthouse/accessibility/}
% \end{itemize}

\end{document}